En esta pregunta usamos la siguiente notación que será utilizada en el
curso. Dados dos strings $u, v$, el string $u \| v$ se define como la
concatenación de $u$ con $v$. Además, dados dos strings binarios $u,
v$ del mismo largo, $u \oplus v$ es el string binario que resulta de
tomar el XOR de $u$ con $v$ componente a componente (lo cual
corresponde a tomar suma en módulo 2). Por ejemplo, si $u = 00100$ y
$v = 11110$, entonces $u \| v = 0010011110$ y $u \oplus v = 11010$.


Sea $n$ un número natural, y suponga que $\{f_k\}_{k \in \{0,1\}^n}$
es una familia de funciones tales que $f_k
: \{0,1\}^n \to \{0,1\}^n$. Note que estas funciones no son
necesariamente biyectivas. A partir de esta familia definimos un
esquema criptográfico $(\textit{Gen}, \textit{Enc}, \textit{Dec})$
sobre los espacios $\mathcal{M} = \mathcal{C} = \mathcal{K}
= \{0,1\}^{2n}$. En particular, la función de encriptación
$\textit{Enc}$ es definida de la siguiente forma considerando un
mensaje $m \in \{0,1\}^{2n}$ tal que $m = m_1 \| m_2$ con $m_1,
m_2 \in \{0,1\}^n$, y una clave $k \in \{0,1\}^{2n}$ tal que $k =
k_1 \| k_2$ con $k_1, k_2 \in \{0,1\}^n$:
\begin{align*}
\textit{Enc}_k(m) \ = \ (m_1 \oplus f_{k_1}(m_2)) \| (m_2 \oplus f_{k_2}(m_1 \oplus f_{k_1}(m_2))). 
\end{align*}
En este problema usted debe responder las siguientes preguntas.
\begin{enumerate}
\item {[2 puntos]} Indique como se define la función $\textit{Dec}$. Es decir, dado una clave $k \in \{0,1\}^{2n}$ y un mensaje cifrado $c \in \{0,1\}^{2n}$, indique cómo se calcula el texto plano $m$ tal que $\textit{Enc}_k(m) = c$.

\item {[4 puntos]} Demuestre que el esquema $(\textit{Gen}, \textit{Enc}, \textit{Dec})$ no es una pseudo-random permutation (PRP) con 2 rondas.  
\end{enumerate}
Note que (b) muestra que $(\textit{Gen}, \textit{Enc}, \textit{Dec})$
no es una PRP independientemente de las propiedades que cumpla la
familia de funciones $\{f_k\}$.

\paragraph{Solución.}                                                                                                                
\begin{enumerate}                                                                                                                   
  \item Sea $c \in \{0,1\}^{2n}$ un mensaje cifrado tal que $c = c_1 \| c_2$ con $c_1,
c_2 \in \{0,1\}^n$, y sea $k \in \{0,1\}^{2n}$ una clave tal que $k =
k_1 \| k_2$ con $k_1, k_2 \in \{0,1\}^n$. Para calcular un mensaje $m$ tal que $\textit{Enc}_k(m) = c$ realizamos los siguientes cálculos:
\begin{align*}
m_1 \ = \ &c_1 \oplus f_{k_1}(c_2 \oplus f_{k_2}(c_1))\\
m_2 \ = \ &c_2 \oplus f_{k_2}(c_1)\\
m \ = \ &m_1 \| m_2
\end{align*}
Para terminar esta parte de la pregunta tenemos que demostrar que $\textit{Enc}_k(m) = c$:
\begin{align*}
\textit{Enc}_k(m) \ = \ &\textit{Enc}_k(m_1 \| m_2)\\
\ = \ &(m_1 \oplus f_{k_1}(m_2)) \| (m_2 \oplus f_{k_2}(m_1 \oplus f_{k_1}(m_2)))\\
\ = \ &(c_1 \oplus f_{k_1}(c_2 \oplus f_{k_2}(c_1)) \oplus f_{k_1}(c_2 \oplus f_{k_2}(c_1))) \|\\
&(c_2 \oplus f_{k_2}(c_1) \oplus f_{k_2}(c_1 \oplus f_{k_1}(c_2 \oplus f_{k_2}(c_1)) \oplus f_{k_1}(c_2 \oplus f_{k_2}(c_1))))\\
\ = \ & c_1 \| (c_2 \oplus f_{k_2}(c_1) \oplus f_{k_2}(c_1))\\
\ = \ & c_1 \| c_2\\
\ = \ & c
\end{align*}

\item Considere la definición de PRP dada en el curso, y suponga que $n \geq 1$. Para demostrar que el esquema definido en esta pregunta no es una PRP con $2$
rondas, usamos un adversario que ejecuta los siguientes pasos:
\begin{itemize}
  \item El adversario entrega $y = 0^{n} \| 0^n$ y $z = 1^n \| 0^n$ al verificador. El adversario recibe como respuesta $f(y)$ y 
    $f(z)$.

  \item Suponiendo que $f(y) = y_1 \| y_2$ con $y_1, y_2 \in \{0,1\}^n$, y
$f(z) = z_1 \| z_2$ con $z_1, z_2 \in \{0,1\}^n$, el adversario verifica si $z_1 = 1^n \oplus y_1$.
    Si esta condición se cumple, entonces el adversario
    indica que $b=0$, sino indica que $b=1$.
\end{itemize}
Necesitamos calcular la probabilidad de que el adversario gane el
juego, lo cual está dado por la siguiente expresión:
\begin{align*}
  \Pr(&\text{Adversario gane el juego}) \ =\\
  &\Pr(\text{Adversario gane el juego} \mid b=0) \cdot \Pr(b=0) \ +\\
  &\hspace{152pt} \Pr(\text{Adversario gane el juego} \mid b=1) \cdot \Pr(b=1) \ =\\
  &\frac{1}{2} \cdot \Pr(\text{Adversario gane el juego} \mid b=0) + \frac{1}{2} \cdot \Pr(\text{Adversario gane el juego} \mid b=1).
\end{align*}
Si $b = 0$, entonces el verificador debe haber encriptado los mensajes
$y$, $z$ con una clave $k \in \{0,1\}^{2n}$ tal que $k = k_1 \| k_2$ con $k_1, k_2 \in \{0,1\}^n$. En particular, se debe tener que:
\begin{align*}
y_1 \ = \ &0^n \oplus f_{k_1}(0^n) \ = \ f_{k_1}(0^n)\\ 
y_2 \ = \ &0^n \oplus f_{k_2}(0^n \oplus f_{k_1}(0^n)) \ = \ f_{k_2}(f_{k_1}(0^n))\\ 
z_1 \ = \ &1^n \oplus f_{k_1}(0^n)\\
z_2 \ = \ &0^n \oplus f_{k_2}(1^n \oplus f_{k_1}(0^n)) \ = \ f_{k_2}(1^n \oplus f_{k_1}(0^n)) 
\end{align*}
Así, se deduce que $z_1 = 1^n \oplus y_1$, y por lo tanto el adversario elige $b = 0$ y gana el juego. Se concluye
que $ \Pr(\text{Adversario gane el juego} \mid b=0) = 1$.  Si $b = 1$,
entonces el verificador escoge al azar una permutación $\pi :
\{0,1\}^{2n} \to \{0,1\}^{2n}$ y responde en el juego con el valor $f(y) =
\pi(y)$ y $f(z) = \pi(z)$. En este caso, tenemos que el adversario pierde el juego si $z_1 = 1^n \oplus y_1$.
Vale decir,
\begin{align*}
  & \Pr(\text{Adversario pierda el juego} \mid b=1)\\
  & \hspace{100pt} = \ \Pr_{\pi}(\pi(y) = y_1 \| y_2 \wedge \pi(z) = z_1 \| z_2 \wedge z_1 = 1^n \oplus y_1) \ =\\
  & \hspace{100pt} = \ \frac{2^{2n} \cdot 2^n \cdot (2^{2n} - 2)!}{(2^{2n})!}\\
  & \hspace{100pt} = \ \frac{2^n}{2^{2n}-1} \ \leq \ \frac{2^n}{2^{2n}-2^n} \ = \ \frac{1}{2^n-1}.
\end{align*}
Entonces se tiene que:
\begin{align*}
  &\Pr(\text{Adversario gane el juego}) \ =\\
  &\hspace{10pt}\frac{1}{2} \cdot \Pr(\text{Adversario gane el juego} \mid b=0) + \frac{1}{2} \cdot \Pr(\text{Adversario gane el juego} \mid b=1) \ \geq\\
  &\hspace{10pt}\frac{1}{2} + \frac{1}{2} \cdot \bigg(1 - \frac{1}{2^n-1}\bigg).
\end{align*}
Por lo tanto, $\Pr(\text{Adversario gane el juego}) \geq \frac{1}{2} + g(n)$, donde $g(n) = \frac{1}{2} \cdot \big(1 - \frac{1}{2^n-1}\big)$ es una función no despreciable. 
Tenemos entonces que  el esquema definido en esta pregunta no es una PRP con $2$
rondas.
\end{enumerate}
